\section*{Chapter \ref{ch:s2:hedging-a-structured-products-book} --- Hedging a Structured-Products Book}

\begin{solution}{exo:s2:hedging-a-structured-products-book:1}
$(100 - 99.12)/100 \times \euro400$ million $\approx$ \euro3.52 million.
\end{solution}

\begin{solution}{exo:s2:hedging-a-structured-products-book:2}
$9.66 / 5 = 1.93$ points; the bank pays half of it on average on the vega it sells: $9.66 \times 1.93 / 2 \approx$ \euro9.33 million.
\end{solution}

\begin{solution}{exo:s2:hedging-a-structured-products-book:3}
It buys \euro228 million of the index as it falls from 100\% to 55\%: the hedge grows from \euro106 million to \euro334 million.
\end{solution}

\begin{solution}{exo:s2:hedging-a-structured-products-book:4}
The investor is short a put that knocks in below the protection level, and the autocall caps the upside it could earn from coupons; the issuer holds
the other side. A put gains value when volatility rises, so the issuer's liability falls: it is long volatility.
\end{solution}

\begin{solution}{exo:s2:hedging-a-structured-products-book:5}
The worst-of pays according to the worse of two indices. The more alike they move, the less likely one of them ends far below the other, so the
note is worth more to the investor; the issuer, short the note, loses when correlation rises.
\end{solution}

\begin{solution}{exo:s2:hedging-a-structured-products-book:6}
Near the barriers the hedge's size changes fastest with the index. Below the level where it peaks the issuer must sell the index as it falls, adding
to the fall, and when many notes share barriers and underlyings, as the BIS described for Korea, the sales are large relative to the market.
\end{solution}

\begin{solution}{exo:s2:hedging-a-structured-products-book:7}
Between 50\% and 55\% of the strike, about 52\% by interpolation (vega $-0.47$ million at 50\%, $+0.76$ million at 55\%). Below it the protection is
likely lost and the note is mostly a holding of the index; more volatility then raises the chance that the index recovers above 60\% by maturity,
which helps the investor and costs the issuer.
\end{solution}

\begin{solution}{exo:s2:hedging-a-structured-products-book:8}
The margin is \euro8.31 million on \euro1 billion; recycling the book's vega alone costs \euro4.67 million at a depth of \euro10 million per point and
\euro9.33 million at \euro5 million. The unhedgeable exposures and the concession needed to pass them on are the economics of the business, not a
detail.
\end{solution}

\begin{solution}{pb:s2:hedging-a-structured-products-book:1}
\begin{enumerate}
\item A sensitivity of a structured-products book to an input vanilla instruments hedge only in part, measured by revaluing the book with it bumped.
\item \euro400 million single on A, \euro300 million single on B, \euro300 million worst-of; five years, annual autocall at 100\%, Phoenix coupons of 7\%
and 9\% at 70\%, 60\% protection; 20\% volatility with skew, 3\% dividends, correlation 0.7.
\item \euro8.31 million: the notes are worth 99.12 and 99.29 per 100.
\item Autocallables are widely sold to Korean retail investors, hedging them can require destabilising trades near barriers, and issuance on Samsung
and SK Hynix rose sharply in the first half of 2026.
\item Vega 3.22, 2.42, 4.02, total 9.66; skew 1.47, 1.10, 1.86, total 4.42; dividends 0.51, 0.38, 0.55, total 1.43; correlation $-1.05$ (\euro{}
millions).
\item It is long the put the investor sold, and a put gains from higher volatility, steeper downside skew and higher dividends.
\item The worst-of's investor gains when the indices move alike; the issuer holds the other side.
\item From \euro3.22 million at the strike to a peak of \euro4.45 million near 80\%, turning negative below about 52\%.
\item Passing on exposures the book cannot hedge with vanillas: selling long-dated volatility and dividends, buying correlation.
\item Vega through long-dated options and variance swaps to hedge funds; dividends through dividend futures and swaps; correlation by buying it
through dispersion trades.
\item At depths of \euro5, 10 and 20 million per point: implied volatility falls 1.93, 0.97 and 0.48 points; the bank pays 9.33, 4.67 and 2.33 million;
a patient buyer gains 18.67, 9.33 and 4.67 million.
\item The buyers of the recycled exposures, paid by the issuers through the concession.
\item Long \euro9.66 million of vega, \euro4.42 million per 0.05 of skew, \euro1.43 million per 0.1 point of dividends, short \euro1.05 million per 0.05 of
correlation; recycling the vega moves long-dated implied volatility by $-0.97$ points at a depth of \euro10 million.
\item Below 55\% the hedge shrinks as the index falls: the bank sells into the fall.
\item Of the same order: \euro4.67 million against a margin of \euro8.31 million at a depth of \euro10 million.
\item They share the depth: the move and the cost grow with the total supply.
\item Chapter~2's dispersion, chapter~3's term-structure and \omterm{def:s2:skew-and-term-structure-relative-value:skew}{skew trades}, and chapter~26's dividend buyers.
\item The vega-recycling buyer's.
\item On quoted long-dated implied volatility with its spreads, entering after issuance waves and holding to reversion, including the years it does not
revert.
\item The exposures its notes leave it with, recycled to whoever will take them at a price.
\end{enumerate}
\end{solution}

\begin{solution}{iq:s2:hedging-a-structured-products-book:1}
Long volatility (mostly long-dated), long downside skew, long dividends, short correlation for worst-of notes, and gap and barrier risk that grows as
the underlying approaches the knock-in levels.
\end{solution}

\begin{solution}{iq:s2:hedging-a-structured-products-book:2}
Revalue the book with the correlation matrix bumped (uniformly and pair by pair), with common random numbers, and express the result as correlation
vega by pair and maturity; check it against the book's dispersion hedges.
\end{solution}

\begin{solution}{iq:s2:hedging-a-structured-products-book:3}
Recompute the exposures at the new level, since vega and delta change fastest there; reduce the delta hedge carefully across the barriers, buy
downside protection if it is not too expensive, and coordinate with the risk limits before the market moves further; expect others with the same
books to trade the same way.
\end{solution}

\begin{solution}{iq:s2:hedging-a-structured-products-book:4}
Large falls through the barriers with volatility and correlation jumping, dividend cuts, and a gap move on an observation date; each with the
liquidity of the hedging markets reduced.
\end{solution}

\begin{solution}{iq:s2:hedging-a-structured-products-book:5}
Share simulated paths across notes with the same underlyings, use adjoint (pathwise) sensitivities instead of bumps, reduce the time grid to the
observation dates where the payoff allows, and run on many cores or GPUs.
\end{solution}

\begin{solution}{iq:s2:hedging-a-structured-products-book:6}
Selling $x$ moves the price by $x/D$; the marginal price of the $v$-th unit is $v/D$, so the cost is $\int_0^V v/D\,dv = V^2/(2D)$. In $n$ parts
with full recovery in between, each costs $(V/n)^2/(2D)$, and $n$ of them cost $V^2/(2nD)$.
\end{solution}
